\section{A polylogarithmic boundary identity at every even order}
\label{sec:polylog-companion}

The new binomial annihilator has an exact polylogarithmic realization.
We use the one-argument, decreasing-index convention
\[
 \operatorname{Li}_{p,q}(z)=\sum_{n>m\ge1}\frac{z^n}{n^p m^q},
 \qquad |z|<1.
\]
Thus only the largest summation index carries the variable. This
convention is essential: a product of two ordinary polylogarithms
would give a different coloring.

\begin{proposition}[The two-tail generating function]
\label{prop:two-tail-polylog}
For integers $p,q\ge2$ and $|z|<1$, define
\[
 \mathcal F_{p,q}(z)=\sum_{n\ge0}
       \bigl(\zeta(p)-H_n^{(p)}\bigr)
       \bigl(\zeta(q)-H_n^{(q)}\bigr)z^n.
\]
Then
\begin{equation}\label{eq:two-tail-polylog}
 \mathcal F_{p,q}(z)=\frac{1}{1-z}\Bigl\{
 \zeta(p)\zeta(q)-\zeta(q)\operatorname{Li}_p(z)
 -\zeta(p)\operatorname{Li}_q(z)
 +\operatorname{Li}_{p,q}(z)+\operatorname{Li}_{q,p}(z)
 +\operatorname{Li}_{p+q}(z)\Bigr\}.
\end{equation}
\end{proposition}
\begin{proof}
Put $T_n=(\zeta(p)-H_n^{(p)})(\zeta(q)-H_n^{(q)})$.
The exact difference is
\[
 T_{n-1}-T_n=n^{-p}(\zeta(q)-H_n^{(q)})
            +n^{-q}(\zeta(p)-H_n^{(p)})+n^{-p-q}.
\]
Multiplication of the generating series by $1-z$ and substitution
of this difference give the result, using
\[
 \sum_{n\ge1}\frac{H_n^{(q)}}{n^p}z^n
   =\operatorname{Li}_{p,q}(z)+\operatorname{Li}_{p+q}(z).
\]
All series converge absolutely and locally uniformly on the open
unit disk, so the rearrangements are justified.
\end{proof}

The generating-function identity is a routine consequence of finite
harmonic summation. The additional content is the particular
all-order cancellation furnished by the determinant theorem.
Write
\[
 \mathcal E=z\frac{d}{dz}+\frac12.
\]
For the polynomial $W_b$ in \eqref{eq:W-polynomial-two}, let
$R_b$ be the nonnegative-power part of the centered asymptotic
Laurent series
\begin{equation}\label{eq:polylog-subtraction}
 W_b(X)\zeta(2,X+\tfrac12)\zeta(2b,X+\tfrac12).
\end{equation}
Only finitely many Bernoulli coefficients enter $R_b$; it is an
explicit even polynomial of degree $2b-2$.

\begin{theorem}[An all-order polylogarithmic boundary evaluation]
\label{thm:polylog-boundary}
For every integer $b\ge2$,
\begin{equation}\label{eq:polylog-boundary}
 \boxed{\displaystyle
 \lim_{z\uparrow1}\left\{
 W_b(\mathcal E)\mathcal F_{2,2b}(z)
       -R_b(\mathcal E)\frac{1}{1-z}\right\}
 =-\frac{2b-1}{4}\bigl(2-B_{2b-2}\bigr).}
\end{equation}
The expression in braces has an absolutely convergent Taylor series
on the closed unit disk, including its value at $z=1$.
\end{theorem}
\begin{proof}
On each compact subset of the open disk, polynomial differentiation
of the locally normally convergent generating series gives
\[
 W_b(\mathcal E)\mathcal F_{2,2b}(z)
       -R_b(\mathcal E)\frac{1}{1-z}
 =\sum_{n\ge0}\left\{
 W_b(n+\tfrac12)A_2(n)A_{2b}(n)-R_b(n+\tfrac12)
 \right\}z^n.
\]
The centered Hurwitz expansion contains only even inverse powers.
After the displayed subtraction its coefficients are $O(n^{-2})$.
The series therefore converges absolutely and uniformly for $|z|\le1$.
Evaluation at $z=1$ is exactly the ordinary convergent identity in
Corollary~\ref{cor:ordinary-two}.
\end{proof}

Equations~\eqref{eq:two-tail-polylog} and
\eqref{eq:polylog-boundary} give a finite prescription involving
classical polylogarithms, depth-two polylogarithms, and their ordinary
derivatives. The subtraction is determined before the limit is taken.
There is consequently no unspecified regularization constant and
no need to assign values to the divergent pieces separately.
